\chapter{Introducción}
\label{cap:introduccion}

Hoy en día los videojuegos han ido mucho más allá del simple entretenimiento para convertirse en un medio audiovisual capaz de transmitir vivencias personales y tratar temas sociales complejos, como la salud mental o la búsqueda de la propia identidad \parencite{bogost2011, navarro2023}. Dentro de este contexto, abordar una crisis de identidad desde la perspectiva de género permite llevar al terreno interactivo el cansancio psicológico que supone intentar encajar en las normas impuestas por la sociedad y lidiar con el rechazo, tanto externo como propio.

Sin embargo, al jugar a muchos títulos que tratan temas parecidos, es habitual encontrar que la historia y la jugabilidad van por caminos separados. Muchas veces toda la carga emocional se cuenta mediante escenas pasivas o textos largos que cortan el ritmo de la partida, mientras que las mecánicas de juego siguen siendo exactamente las mismas de principio a fin sin importar cómo se sienta el protagonista \parencite{finestrat2018}. Por otro lado, en aquellos juegos que sí intentan cambiar sus reglas de forma oculta según cómo jugamos, suele aparecer otro problema: al modificar los controles sin avisar, el jugador siente que le quitan el mando de las manos y acaba frustrándose.

Para buscar una solución a ambos problemas en este \ac{TFG}, realizado en la \ac{EPS}, se plantea el diseño y desarrollo de un prototipo funcional en \ac{2D} visto en primera persona con sensación de profundidad (2.5D). En él, un sistema en segundo plano —el Director de Alucinación— observa cómo reacciona el usuario ante los obstáculos a través de tres dimensiones continuas (ímpetu, cautela y templanza) y, pidiendo siempre permiso previo mediante una mecánica de consentimiento explícito, adapta en tiempo real las capacidades del protagonista, la respuesta de los mundos y la interfaz gráfica de usuario (\ac{UI}) a través de distintos «Mecanismos de Defensa».

Esta memoria se ha organizado siguiendo las pautas del Libro de Estilo de la Escuela Politécnica Superior de la Universidad de Alicante \parencite{eps2015} y tomando como guía la estructura habitual de los trabajos del Grado en Ingeniería Multimedia \parencite{finestrat2018}:

\begin{itemize}
    \item \textbf{Capítulo 1. Introducción:} presenta el tema del proyecto, explica el problema detectado al unir narrativa y mecánicas de juego y resume la propuesta planteada.
    \item \textbf{Capítulo 2. Marco teórico o Estado del arte:} repasa la base teórica sobre narrativa digital, narrativa ambiental, armonía y disonancia ludonarrativa, salud mental desde la perspectiva de género y sistemas adaptativos, además de analizar varios videojuegos que han servido de referencia.
    \item \textbf{Capítulo 3. Objetivos:} detalla el objetivo general del trabajo y desglosa los objetivos específicos que guían tanto el diseño como la programación del prototipo.
    \item \textbf{Capítulo 4. Metodología:} explica paso a paso las cuatro fases en las que se divide el proyecto, desde la documentación inicial y la consulta psicológica hasta la justificación del motor Unity y el resto de herramientas.
    \item \textbf{Capítulo 5. Diseño del Videojuego (GDD) y Dirección Artística:} desarrolla todo el diseño del juego, incluyendo el giro argumental del falso apocalipsis zombi, el funcionamiento de los tres mundos, las reglas del Director de Alucinación, el sistema de consentimiento y el estilo visual en primera persona con \textit{billboards} pintados a mano.
    \item \textbf{Capítulo 6. Arquitectura Técnica e Implementación:} documenta la estructura técnica en Unity, la programación de los sistemas adaptativos y el montaje de la demo jugable.
    \item \textbf{Capítulo 7. Conclusiones:} resume qué objetivos se han cumplido con el prototipo y qué mejoras o ampliaciones quedan abiertas para continuar con el juego en el futuro tras el \ac{TFG}.
\end{itemize}